\chapter{Wstęp}



\section{Cel pracy}
Celem niniejszej pracy inżynierskiej jest stworzenie systemu sterującego, który będzie w stanie przeprowadzić stabilny pionowy lot rakiety. System sterujący rakietą oparty jest o samouczący algorytm sztucznej inteligencji. Aby osiągnąć główny cel pracy konieczne jest stworzenie środowiska, w którym algorytm sterujący będzie trenowany. Stworzenie silnika fizycznego umożliwiającego symulację lotu rakiety, oraz ingerencję w jej lot poprzez zmianę parametrów układu sterującego rakiety jest pośrednim celem niniejszej pracy. 

\section{Ogólny zarys pracy}

Poniższa praca inżynierska składa się z 7 rodziałów... [TO JEST DO DOPISANIA NA KONIEC]

\section{Użyte technologie}
Projekt tworzony w ramach niniejszej pracy inżynierskiej składa się z dwóch głównych modułów, stanowiących rozwiązania spełniające odpowiednio główny i pośredni cel pracy. 
\\Następujące narzędzia i technologie zostały wykorzystane do stworzenia oraz integracji obydwu modułów projektu:

\begin{itemize}
    \item \textbf{Git oraz Github} - narzędzia wykorzystane do kontroli wersji projektu.
    \item inne narzędzia wykorzystane do stworzenia całości projektu
\end{itemize}
Poniżej opisane zostały poszczególne narzędzia i technologie wykorzystane w każdym z modułów.

\subsection{Silnik fizyczny symulacji lotów rakietowych}
Do stworzenia silnika fizycznego symulacji lotów rakietowych wykorzystane zostały następujące technologie:

\begin{itemize}
    \item \textbf{Python 3.13} - język programowania 
    \item zestaw bibliotek  
    \begin{itemize}
        \item bibioteka1....
    \end{itemize}
\end{itemize}
    
\subsection{Algorytm sterujący rakietą}
Do stworzenia algorytmu sterującego rakietą wykorzystane zostały następujące technologie: 

\begin{itemize}
    \item gówno
    \item dupa
\end{itemize}



\section{Przegląd literatury}

Przegląd literatury podzielony został na dwie części dotyczące dwóch modułów projektu. W pierwszej części opisana została literatura dotycząca tworzenia silnika fizycznego symulacji lotów rakietowych, a druga część zawiera opis literatury odnoszącej się do tworzenia algorytmu sterującego rakietą.

\subsection{Silnik fizyczny symulacji lotów rakietowych}

Aby dokładnie zamodelować fizykę lotów rakietowych niezbędna jest wiedza na temat sił działających na obiekt poruszający się w atmosferze, w polu grawitacyjnym.
\subsubsection{Dostępne rozwiązania}
\subsubsection{Dynamika}
\subsubsection{Aerodynamika}


\subsection{Algorytm sterujący rakietą}